% !PRINTED 58
\weberpage{58}

\begin{center}
\S.\,15.

\gesperrt{Indices for a Power of 2 as Modulus}.
\end{center}

A similar theorem must now be established for the powers $2^\lambda$ of the prime 2, which, as already remarked above, behaves differently.

If $\lambda = 1$, then every number not divisible by 2 is $n \equiv 1 \pmod{2}$. If $\lambda = 2$, then $-1$ is to be regarded as a primitive root of 4, for every odd number $n$ satisfies one of the congruences

\[
n \equiv (-1)^\alpha \pmod{4},
\]

in which $\alpha = 0$ or $= 1$.

But already for the modulus 8 no primitive root exists any longer, that is, no number by whose powers all classes of odd numbers with respect to the modulus 8 can be represented. For if $g$ is any odd number, then among the powers of $g$ at most the two $1, g$ are distinct with respect to the modulus 8, because always $g^2 \equiv 1 \pmod{8}$. If therefore one takes $g$ not congruent to 1 and not congruent to $-1 \pmod{8}$, thus $g = 3$ or $= 5$, then every odd number is congruent to one of the four numbers

\[
(-1)^\alpha g^\beta \qquad \alpha = 0, 1, \quad \beta = 0, 1
\]

with respect to the modulus 8.

The number of classes of odd numbers with respect to the modulus $2^\lambda$ is $2^{\lambda-1}$. Now, however, for every odd number $g$, provided $\lambda > 2$,

\weq{1}{g^{2^{\lambda-2}} \equiv 1 \pmod{2^\lambda}.}

For if we assume (1) to be correct and accordingly set

\[
g^{2^{\lambda-2}} = 1 + h\, 2^\lambda,
\]

and raise to the square, it follows that

\[
g^{2^{\lambda-1}} \equiv 1 \pmod{2^{\lambda+1}}.
\]

If therefore the formula (1) is correct for $\lambda$, then it is so also for $\lambda + 1$, and since it holds for $\lambda = 3$, it holds generally. From this it follows that among the powers of an odd number $g$ at most $2^{\lambda-2}$ distinct ones with respect to the modulus $2^\lambda$ can occur.

On the other hand, however, it follows easily for $g = 5$ that $5^{2^{\lambda-2}}$ is the lowest power that becomes congruent to unity with respect to the modulus $2^\lambda$. For if $5^e$ is the lowest power that is congruent to
%CONT>>
